\chapter{Bilder, Tabellen, Formeln}%
%
%
\section{Label und Referenzen}%
\label{sec: labandref}%
%
%
\section{Tabellen}%
\label{sec: tables}%
%
%
\section{Bilder}%
Bilder müssen wie auch Tabellen zum sinnvollen Beschriften und Referenzieren in eine Umgebung geschrieben werden. Nur so kann man Bilder referenzieren (siehe \ref{sec: labandref} und sie wie auch Tabellen benennen (siehe \ref{sec: tables}).\\%

\begin{figure}[h]%
    \includegraphics[width=0.25\textwidth]{stura-logo-htw.pdf}%
    \caption{Logo des STURA der HTW Dresden}%
    \label{fig: stura-logo}%
\end{figure}%
dann noch Bilder:
Diese müssen in der Diplomarbeit - egal wie winzig sie sind - eine Bildunterschrift haben und ebenso in einer eigenen Zeile stehen!!! Keinen Umlaufenden Text verwenden.
%
%
\section{Positionieren von Bildern und Tabellen}%Vielleicht auch von Formeln? uff, doch noch viel zu tun.
\section{Formeln}%
Die wichtigsten Dinge für Mathematik im Dokument sind folgende:
aus\\
\begin{tabular}[h]{cc}
\hline 
 a=b  & \textbackslash begin \{equation\} a=b \textbackslash end \{equation\} \\ 
\hline 
\end{tabular} 

\begin{lstlisting}
\begin{equation}
a=b
\end{equation}
\end{lstlisting} 
wird
\begin{equation}
a=b
\end{equation}
aus 
\begin{lstlisting}
\begin{equation*}
a=b
\end{equation*}
\end{lstlisting}
wird
\begin{equation*}
a=b
\end{equation*}
ohne Nummerierung

Mathe im Text schreibt man mit dem Dollarzeichen
\begin{lstlisting}
$E = mc^2$
\end{lstlisting}
und man erhält somit $E=mc^2$.

Damit sollten die wichtigsten Arten für Mathematik in der Diplomarbeit abgedeckt sein.
%
%
\section{Literatur}%
\label{sec: literatur}
Die Literatur schreibt man in die Datei \emph{literatur.bib}, welche relativ zu eurem Template im Verzeichnis \emph{/frontbackmatter/} liegt. Hier kommt euch \emph{TexMaker} zu Hilfe, denn unter \emph{Bibliography -> biblatex} findet man Verschiedene Arten von Literatur, die man mit einem Klick hinzufügen kann.%

Dabei sollte man beachten, dass zwischen den einzelnen Quellen unbedingt eine \textsc{Leerzeile} sein sollte, da es sonst eventuell zu Problemen beim Referenzieren kommt.%

Im ersten Feld gibt man seinen \emph{Key} für diese Quelle an. Der \emph{Key} funktioniert später beim zitieren wie ein \emph{Label}. Man kann damit auf den entsprechenden Eintrag verweisen. Alle anderen Felder sind größtenteils selbsterklärend und können sonst in der Dokumentation für Biblatex nachgelesen werden.%

Anschließend kann man mit der Option \emph{Bibliography -> Clean} alle optionalen Felder automatisch entfernen lassen, sofern die Buchstaben \emph{OPT} am Anfang der Zeile nicht entfernt wurden.%
%
%
\section{Zitieren}%
\label{sec: zitieren}
Nachdem man seine Quellen dem Literaturverzeichnis hinzugefügt hat (\ref{sec: literatur})und diese mit \emph{Keys} versehen hat versehen hat, kann man diese anschließend Zitieren.%
Für indirekte Zitate muss man nur die Quelle am Ende des Absatzes/Satzes angeben. Wie das geht, sieht man in \ref{tab: indirekteszitat}. Bemerkungen kann man einfügen, muss man aber nicht. Als \emph{Notiz nachher} empfiehlt es sich, z.B. eine Seitenzahl anzugeben wie hier am Ende des Satzes zu sehen. \cite[][S.22 ff]{einstein}
\begin{table}[h]%
\begin{tabu} to \textwidth {X[c]X[c]}%
\toprule%
Hier steht ein indirektes Zitat.\textbackslash cite[Notiz vorher][Notiz nachher]\{key\} & Hier steht ein indirektes Zitat.\cite[Notiz vorher][Notiz nachher]{einstein}\\%
\bottomrule%
\end{tabu}%
\caption{Indirektes Zitat}%
\label{tab: indirekteszitat}%
\end{table}%

Manchmal möchte man unbedingt manuell Anführungszeichen setzen, wie in \ref{tab: quotationmarks} zu sehen. Man beachte, dass man immer ein Zitat öffnen und schließen muss. Auf ein \guillemotright \space sollte immer später ein \guillemotleft \space folgen, da es sonst zu Fehlern beim Kompilieren kommt.\\%
\begin{table}[h]%
\begin{tabu} to \textwidth {X[c]X[c]X[c]}%
\toprule%
\textbf{Tastenkombination} & \textbf{Latex Code} & \textbf{Befehl}\\%
\midrule%
SHIFT+F1 & \guillemotright & Anführungszeichen doppelt öffnen\\ %
SHIFT+F2 & \guillemotleft & Anführungszeichen doppelt schließen\\%
SHIFT+F3 & \guilsinglright & Anführungszeichen einfach öffnen\\%
SHIFT+F4 & \guilsinglleft & Anführungszeichen einfach schließen\\%
\bottomrule%
\end{tabu}%
\caption{Manuelle Anführungszeichen}%
\label{tab: quotationmarks}%
\end{table}%

Für wörtliche Zitate nutzt man den Befehl aus \ref{tab: woertlicheszitat}. Mit diesem Befehl werden Zitate entweder in Anführungszeichen gesetzt oder automatisch in Blöcke geschrieben, sobald sie 3 Zeilen überschreiten. Möchte man Anführungszeichen trotz Blockschreibweise haben, so fügt man die Anführungszeichen aus \ref{tab: quotationmarks} hinzu.\\%
\begin{table}[h]%
\begin{tabu} to \textwidth {X[c]X[c]}%
\toprule%
\textbackslash blockcquote[vorher][nachher]\{key\}\{Hier steht ein wörtliches Zitat.\} & \blockcquote[vorher][nachher]{einstein}{Hier steht ein wörtliches Zitat.}\\%
\bottomrule%
\end{tabu}%
\caption{Wörtliches Zitat}%
\label{tab: woertlicheszitat}%
\end{table}%

Die Auslassungspunkte, die man häufig beim Zitieren verwendet, entstehen folgendermaßen (siehe \ref{tab: auslassungspunkte}):\\%
\begin{table}[h]%
\begin{tabu} to \textwidth {X[c]X[c]}%
\toprule%
\textbackslash textelp\{\} & \textelp{} \\%
\bottomrule%
\end{tabu}%
\caption{Auslassungspunkte}%
\label{tab: auslassungspunkte}%
\end{table}%


Für Zitate in anderen Sprachen schaut man in die Dokumentation von dem Paket \emph{csquotes}.%

Damit die Literaturreferenzen im Dokument erkannt werden, muss man vor dem Kompilieren den selbst erstellten \emph{User Command} für \emph{biblatex} starten und anschließend die PDF mindestens zwei mal Kompilieren, damit das Ergebnis vollständig sichtbar wird.%
%
%
\section{Fußnoten}%
Fußnoten setzt man wie in \ref{tab: fussnoten} zu sehen.\\%
\begin{table}[h]%
\begin{tabu} to \textwidth {X[c]X[c]}%
\toprule%
Wort \textbackslash footnote\{Fußnotentext\} & Wort \footnote{Fußnotentext} \\%
\bottomrule%
\end{tabu}%
\caption{Fußnoten}%
\label{tab: fussnoten}%
\end{table}%
%
%
\section{Testsection}

\gls{latex}
\Gls{formula}


\newpage
